\subsection{Global Sequentially Consistent Order}
All (:ref:base.instructions.AFENCE:) operations and all \texttt{SEQCST} atomic operations participate in one global total order that is
consistent with each (:term:base.terms.logical_processor:)\textquotesingle{}s ordered-before relation and with per-location coherence. A failed
\texttt{CMPXCHG.SEQCST} participates in that order as an SC load even though it performs no write.

A naturally aligned tear-free scalar load or store also participates as an SC operation when it is the only memory
access between its immediately surrounding (:ref:base.instructions.AFENCE:) operations in memory-event order. Thus
\texttt{AFENCE; access; AFENCE} is the architectural sequence for an SC scalar load or store. The access and both
fences occupy their required positions in the same global order.

Acquire and release operations may be strengthened with (:ref:base.instructions.AFENCE:); the resulting instruction sequence has the
union of the ordering constraints imposed by its operations.
Leaf (:ref:base.memory.translation.PTE.AM:) and (:ref:base.memory.translation.PTE.CP:) are orthogonal. An MMIO access retains its non-speculative,
operation-class, alignment, and ordering rules under every (:ref:base.memory.translation.PTE.CP:) value. MMIO
accesses from one (:term:base.terms.logical_processor:) are observed in program order relative to
other MMIO accesses from that (:term:base.terms.logical_processor:). Normal-memory and MMIO
accesses retain the baseline weak ordering between classes; (:ref:base.instructions.AFENCE:)
orders earlier accesses before later accesses across that boundary.
